\section{Discussion}

The case studies suggest that recursive recovery is most useful when it is treated as a rule-governed change of scientific state, not as persistence. Reopening after a donor closure without a new residual would be search inflation. Refusing to reopen after a representation or access boundary changes would treat an old negative as globally permanent. Typed state and prospectively declared reopen conditions provide the bridge between those extremes.

The claim-preserving recovery theorem sharpens this distinction. If a failed claim identity has load-bearing predicates $P_1,\ldots,P_k$ and the admissible repair language declares which causal ancestors each action can affect, then a same-identity repair must cover at least one ancestor of every failed predicate. An uncovered failed predicate is therefore a no-repair certificate for that action language. The minimum weighted ancestor cover supplies a lower bound on repair effort, and an action whose causal-effect set is contained in that of a no-more-expensive action can be removed from at least one minimum cover. These are necessary structural conditions, not a guarantee that a covered repair will succeed.

The theorem also gives a precise place to draw the successor boundary. Changing the scientific question, population, estimand, protocol semantics, primary metric, threshold, protected corpus or terminal semantics after failure may be a productive research move, but it creates a different claim identity. A successful changed-identity run should therefore be preserved as a successor rather than narrated as a repair of the old failure.

The R6 history illustrates a second distinction: saturation and explanation are not the same. Several exact negatives can establish that a family adds no value on a scope, but a later dominance theorem or counterexample is needed to explain why. A receipts-first archive makes that transition possible because the negative sequence remains available to be compared with the later structural account.

Third, the method separates provenance from authority. Reproducible evidence can still be badly designed, use a weak donor, or support only a narrow inference. Conversely, an honest \textsc{CannotCheck} can be methodologically correct even though it offers no scientific answer. The independent replay results strengthen attribution of the registered computations; they do not turn one programme into evidence of cross-domain effectiveness.

The unresolved population question is comparative: does this discipline improve the reliability, efficiency, or calibration of research programmes relative to other workflows? The current evidence does not answer it. Doing so would require a prospectively defined multi-programme evaluation with matched tasks, matched action/information budgets and independently adjudicated outcomes.